\chapter{Lung conditions}

\section*{Definition and Overview}
Lung conditions encompass a diverse spectrum of pathological disorders affecting the respiratory system, specifically targeting the airways, parenchyma, alveoli, and pulmonary vasculature. These conditions impair the primary physiological function of the lungs: gas exchange, which involves oxygenating blood and removing carbon dioxide. Lung diseases are broadly categorised into obstructive airway diseases (e.g., asthma, chronic obstructive pulmonary disease), restrictive lung diseases (e.g., interstitial lung diseases), infectious diseases (e.g., pneumonia, tuberculosis), pulmonary vascular diseases (e.g., pulmonary embolism, pulmonary hypertension), and neoplastic conditions (e.g., lung cancer). Together, they represent a massive global healthcare burden, causing chronic disability, reducing quality of life, and serving as leading causes of mortality worldwide.

\section*{Epidemiology}
Respiratory diseases are major drivers of morbidity and mortality in Australia and globally. According to the Australian Institute of Health and Welfare (AIHW), chronic respiratory conditions affect approximately 30\% of the Australian population. Asthma is highly prevalent, affecting 11\% of Australians (approximately 2.7 million people), with one of the highest rates of childhood asthma in the world. Chronic Obstructive Pulmonary Disease (COPD) is the fifth leading cause of death in Australia, primarily affecting individuals over the age of 40 who have a history of smoking. Lung cancer remains the leading cause of cancer-related mortality in Australia, accounting for nearly 19\% of all cancer deaths. Globally, acute lower respiratory tract infections, particularly pneumonia, represent a leading cause of mortality in children under the age of five and in adults over the age of 65.

\section*{Aetiology and Pathophysiology}
The development of lung conditions is multifactorial, arising from a combination of environmental exposures, genetic susceptibilities, infectious agents, and lifestyle factors.
\begin{itemize}
    \item \textbf{Tobacco smoke:} The single most significant risk factor for COPD, lung cancer, and several interstitial lung diseases. Cigarette smoke induces chronic airway inflammation, cilia destruction, alveolar wall damage (emphysema), and oncogenic DNA mutations.
    \item \textbf{Environmental and occupational exposures:} Inhalation of hazardous dusts, such as crystalline silica (causing silicosis), asbestos fibres (causing asbestosis and mesothelioma), and coal dust (causing coal workers' pneumoconiosis). Air pollution and particulate matter also exacerbate chronic airway diseases.
    \item \textbf{Infectious pathogens:} Viruses (such as influenza, respiratory syncytial virus [RSV], and SARS-CoV-2), bacteria (such as \textit{Streptococcus pneumoniae} and \textit{Mycobacterium tuberculosis}), and fungi cause acute alveolar and parenchymal inflammation, leading to exudate accumulation and consolidation.
    \item \textbf{Allergen exposure:} Exposure to airborne allergens (such as dust mites, pollens, mould spores, and animal dander) triggers immunoglobulin E (IgE)-mediated mast cell activation and chronic eosinophilic inflammation in susceptible airways, leading to asthma.
    \item \textbf{Genetic determinants:} Mutations in the cystic fibrosis transmembrane conductance regulator (CFTR) gene cause cystic fibrosis, leading to thick, viscous mucus secretions. Alpha-1 antitrypsin deficiency is a genetic disorder that accelerates the development of early-onset emphysema.
\end{itemize}

\section*{Clinical Features}
While clinical presentations vary depending on the specific respiratory disorder, several cardinal symptoms are common to most lung conditions.
\begin{itemize}
    \item \textbf{Dyspnoea:} Shortness of breath, which may occur during physical exertion (characteristic of chronic diseases) or at rest (typical of acute events like pneumothorax, pulmonary embolism, or acute asthma attacks).
    \item \textbf{Chronic cough:} A cough lasting longer than eight weeks. It may be productive of sputum (common in COPD, bronchiectasis, and chronic bronchitis) or dry and hacking (typical of interstitial lung disease or asthma).
    \item \textbf{Wheezing and stridor:} High-pitched musical noises. Wheezing indicates lower airway narrowing (asthma, COPD), while stridor indicates upper airway obstruction.
    \item \textbf{Haemoptysis:} Coughing up blood, which is a red-flag symptom requiring urgent investigation to rule out lung cancer, tuberculosis, pulmonary embolism, or severe bronchiectasis.
    \item \textbf{Chest pain:} Pleuritic chest pain (sharp pain worsened by deep breathing or coughing) is common in pneumonia, pleurisy, pneumothorax, and pulmonary embolism.
    \item \textbf{Systemic symptoms:} Unexplained weight loss, night sweats, fatigue, and persistent low-grade fevers can indicate chronic infections like tuberculosis, or malignancies like lung cancer.
\end{itemize}

\section*{Differential Diagnosis}
Because cardiorespiratory symptoms overlap significantly, clinicians must systematically differentiate lung conditions from other system disorders.
\begin{itemize}
    \item \textbf{Congestive heart failure (CHF):} Presents with exertional dyspnoea, orthopnoea (shortness of breath when lying flat), and paroxysmal nocturnal dyspnoea, distinguished by crackles in the lungs, elevated brain natriuretic peptide (BNP), and cardiomegaly on imaging.
    \item \textbf{Gastro-oesophageal reflux disease (GORD):} A frequent non-pulmonary cause of chronic cough, which may be worse when supine and is accompanied by retrosternal burning.
    \item \textbf{Deconditioning and obesity:} Can cause significant exertional dyspnoea, but lung function testing (spirometry) and chest imaging are normal.
    \item \textbf{Vocal cord dysfunction (inducible laryngeal obstruction):} Can mimic asthma with sudden-onset wheezing and dyspnoea, but does not respond to bronchodilators and shows characteristic upper airway flattening on flow-volume loops.
\end{itemize}

\section*{When to Seek Medical Care}
Patients with chronic lung conditions should monitor their baseline symptoms. Acute respiratory distress requires immediate emergency intervention.

\textbf{Call 000 (or local emergency services) for:}
\begin{itemize}
    \item Sudden, severe, and worsening shortness of breath.
    \item Inability to speak in full sentences or complete thoughts without pausing for breath.
    \item Sharp, crushing, or squeezing chest pain, especially if it radiates to the arm, neck, or jaw.
    \item Blue or grey discolouration of the lips, face, or fingernails (cyanosis).
    \item Coughing up moderate-to-large amounts of blood.
    \item Sudden onset of severe confusion, extreme drowsiness, or difficulty waking up.
\end{itemize}

\section*{Diagnosis and Investigations}
A structured approach combining physiological testing, imaging, and laboratory analysis is required to diagnose respiratory disorders.
\begin{itemize}
    \item \textbf{Spirometry:} The essential test of lung function. It measures the volume of air an individual can exhale forcefully. An FEV1/FVC ratio of less than 0.70 post-bronchodilator confirms airway obstruction (COPD, asthma). Restrictive disorders show reduced forced vital capacity (FVC) with a normal or elevated ratio.
    \item \textbf{Chest X-ray:} The initial imaging test, capable of identifying pneumonia, pleural effusions, pneumothorax, large masses, and severe emphysema.
    \item \textbf{High-resolution computed tomography (HRCT):} Provides detailed cross-sectional views of the parenchyma, crucial for diagnosing interstitial lung disease, bronchiectasis, and staging lung nodules.
    \item \textbf{Diffusing capacity of the lung for carbon monoxide (DLCO):} Measures the transfer of gas from the alveoli to the red blood cells, indicating the integrity of the alveolar-capillary membrane.
    \item \textbf{Sputum culture and cytology:} Evaluates sputum for bacterial, mycobacterial, or fungal infections, and screens for malignant cells.
    \item \textbf{Arterial blood gas (ABG) analysis:} Measures arterial oxygen (PaO2), carbon dioxide (PaCO2), and pH, indicating the severity of acute respiratory failure and acid-base disturbances.
\end{itemize}

\section*{Management}
Management strategies depend on the specific diagnosis, focusing on symptom control, slowing disease progression, and emergency rescue.
\begin{itemize}
    \item \textbf{Pharmacological therapies:} Includes short-acting beta-2 agonists (SABAs) for rapid bronchodilation; long-acting beta-2 agonists (LABAs) and long-acting muscarinic antagonists (LAMAs) for maintenance control in COPD and asthma; inhaled corticosteroids (ICS) to reduce airway inflammation; and antifibrotic agents (e.g., nintedanib) to slow decline in pulmonary fibrosis.
    \item \textbf{Smoking cessation:} The single most effective intervention to stop or slow the decline of lung function in COPD and reduce the risk of lung cancer. Supported by nicotine replacement therapy (NRT) and behavioural support.
    \item \textbf{Pulmonary rehabilitation:} A multidisciplinary programme combining structured exercise training, education, and breathing techniques to improve exercise tolerance and reduce dyspnoea.
    \item \textbf{Oxygen therapy:} Long-term oxygen therapy (LTOT) is indicated for patients with chronic, severe resting hypoxaemia to reduce mortality and strain on the right heart.
    \item \textbf{Surgical and advanced interventions:} Lung volume reduction surgery (LVRS) or endobronchial valve placement for severe emphysema, and lung transplantation for end-stage lung disease.
\end{itemize}

\section*{Complications}
Chronic or severe lung conditions can lead to secondary systemic and organ-specific complications.
\begin{itemize}
    \item \textbf{Respiratory failure:} Inability of the respiratory system to maintain adequate oxygenation (Type I) or carbon dioxide clearance (Type II), requiring mechanical ventilation.
    \item \textbf{Pulmonary hypertension and cor pulmonale:} Chronic hypoxemia causes pulmonary vasoconstriction, increasing pulmonary artery pressure and leading to right ventricular hypertrophy and right-sided heart failure.
    \item \textbf{Pneumothorax:} Rupture of an emphysematous bulla or subpleural bleb, causing air to enter the pleural space and collapse the lung.
    \item \textbf{Recurrent respiratory infections:} Chronic structural damage (e.g., in bronchiectasis or COPD) predisposes patients to repeated bacterial infections, accelerating lung function decline.
\end{itemize}

\section*{Prognosis and Outcomes}
The prognosis of lung conditions varies widely. Asthma generally carries an excellent prognosis with normal life expectancy, provided the patient adheres to controller medications and has an active Asthma Action Plan. COPD is a progressive, irreversible disease, but early intervention and smoking cessation significantly improve survival and quality of life. Interstitial lung diseases, such as idiopathic pulmonary fibrosis (IPF), have a poorer prognosis, with a median survival of three to five years without treatment or lung transplantation. Lung cancer prognosis depends heavily on the stage at diagnosis; early-stage, resectable disease offers a chance of cure, whereas advanced metastatic disease has low long-term survival rates.

\section*{Prevention and Self-Care}
Preventative actions can reduce the incidence of lung conditions and limit exacerbations of chronic disease.
\begin{itemize}
    \item \textbf{Avoid tobacco and vaping:} Quit smoking and avoid exposure to second-hand smoke to protect lung parenchyma and airways from chronic irritation and oncogens.
    \item \textbf{Vaccination:} Stay up to date with annual influenza vaccines, pneumococcal vaccines, and COVID-19 vaccinations to prevent severe respiratory infections.
    \item \textbf{Occupational protection:} Use appropriate personal protective equipment (PPE), such as fitted respirators, when working in environments with silica, asbestos, coal, or chemical fumes.
    \item \textbf{Follow action plans:} Regularly review and follow a written Asthma or COPD Action Plan created with your general practitioner, adjusting medications during exacerbations.
    \item \textbf{Maintain physical activity:} Exercise regularly to build cardiovascular fitness and improve the efficiency of oxygen utilization in body tissues.
\end{itemize}

\section*{Sources and Further Reading}
The following reputable organisations provide comprehensive resources and clinical guidelines on lung health and respiratory diseases:
\begin{itemize}
    \item Lung Foundation Australia. \textit{Support and resources for lung health.} https://lungfoundation.com.au/
    \item National Asthma Council Australia. \textit{Asthma management and resources.} https://www.nationalasthma.org.au/
    \item Healthdirect Australia. \textit{Lung conditions.} https://www.healthdirect.gov.au/lung-conditions
    \item NHS. \textit{Lung conditions and respiratory diseases.} https://www.nhs.uk/conditions/
    \item World Health Organization (WHO). \textit{Chronic respiratory diseases.} https://www.who.int/health-topics/chronic-respiratory-diseases
\end{itemize}
